\chapter{Protocol Variants and Design Choices}
\label{chap:variants}

\section{Motivation for protocol variants}

The inherited implementation supported the main optimistic and dispute flows, but it did so
through contracts that were close to a ``one-size-fits-all'' design. This was convenient for
experimentation, because a single contract family could execute many scenarios. It was less
suitable for gas optimization. A contract that supports every role configuration must keep
state variables, checks, branches, and deployment bytecode that are irrelevant for many
exchanges. The central design question of the project therefore became: which variants are
protocolically meaningful, and which of them justify a separate implementation path?

The answer is not simply the Cartesian product of all possible role equalities. Some choices
are known at precontract time, before the optimistic contract is created. Other choices are
only known if the optimistic path fails and the parties enter the dispute path. The final
implementation separates these two moments. Precontract-time variants select the optimistic
contract to deploy or clone. Dispute-time variants select the appropriate dispute contract
when the dispute sponsors actually appear.

Table~\ref{tab:variant-classification} summarizes the classification used in the
implementation. The table is also the convention used in the rest of the report.

\begin{table}[htbp]
  \centering
  \begin{tabularx}{\linewidth}{>{\raggedright\arraybackslash}p{0.22\linewidth}>{\raggedright\arraybackslash}p{0.25\linewidth}>{\raggedright\arraybackslash}X}
    \toprule
    Variant & Decision time & Main effect \\
    \midrule
    Normal sponsored SOX &
    Precontract / deployment &
    Baseline model with an optimistic sponsor \(S\) and potentially distinct dispute
    sponsors \(S_B\) and \(S_V\). \\
    \texttt{no\_S\_deposit} / simSOX &
    Precontract / deployment &
    Removes the optimistic sponsor deposit and the corresponding reimbursement logic from
    the optimistic phase. \\
    \(S=B\) &
    Precontract / deployment &
    Makes the buyer the optimistic sponsor and allows the implementation to fuse the buyer
    payment with the initial creation step. \\
    \(S=V\) &
    Precontract / deployment &
    Makes the vendor the optimistic sponsor. The buyer payment remains a later optimistic
    action in the current implementation. \\
    \(S_B=B\) &
    After buyer dissatisfaction &
    Makes the buyer its own dispute sponsor and makes the Step 4--5 fusion trivial. \\
    \(S_V=V\) &
    After buyer dissatisfaction &
    Makes the vendor its own dispute sponsor. Together with \(S_B=B\), it removes the need
    for a sponsor-rotation repetition in Step 9. \\
    Hardcoded \(\mathsf{SHA256}\) &
    Precontract / deployment &
    Replaces the generic circuit-opening check for \(h_{\mathit{circuit}}\) in Step 8 by an
    on-chain reconstruction of the expected \(\mathsf{SHA256}\) circuit gate. \\
    \bottomrule
  \end{tabularx}
  \caption{Protocol variants retained for implementation and measurement.}
  \label{tab:variant-classification}
\end{table}

The guiding principle is that a variant is worth implementing separately only if it changes
the cost, risk, or latency of a protocol path. Exotic role equalities that do not change the
execution path are still valid theoretical cases, but they do not necessarily deserve their
own contract. This choice keeps the implementation close to the supervisor feedback:
optimize the cases that matter for measurement, and avoid paying gas for branches that are
not used in the selected mode.

\section{Normal sponsored SOX}

The normal sponsored variant is the reference point. It keeps the full role structure of the
SOX protocol: vendor \(V\), buyer \(B\), optimistic sponsor \(S\), buyer-side dispute sponsor
\(S_B\), and vendor-side dispute sponsor \(S_V\). In this mode, the optimistic sponsor creates
the exchange contract and deposits the funds needed by the optimistic mechanism. The buyer
then pays, the vendor releases the key, and the exchange either completes or moves toward a
dispute.

In implementation terms, the final codebase keeps a normal clone implementation for
factory-based measurements, while the interactive application keeps the direct normal
account path for compatibility with the existing ERC-4337 flow. In both cases, the normal
path starts in state \texttt{WaitPayment}. This means that Step 1 and Step 2 remain distinct
in the normal sponsored path: the sponsor creates and funds the exchange first, and the
buyer later sends the payment.

Keeping this variant is important for two reasons. First, it is the closest implementation
of the sponsored protocol as inherited from the original codebase. Second, it provides a
fair baseline against which all simplified modes can be compared. A simplified mode is only
useful if it is clear which part of the normal sponsored logic it removes.

The normal path also remains the fallback for disputes with external sponsors. When the
buyer-side and vendor-side dispute sponsors are not both self-sponsored, the optimistic clone
deploys the normal dispute contract. This preserves the generic Step 9 behavior, including
the possibility of repeated Step 7--8 executions when the protocol logic requires a new
sponsor configuration.

\section[No S-deposit and simSOX]{No \(S\)-deposit and simSOX}

The \texttt{no\_S\_deposit} mode is the simplified optimistic mode used in the report under
the name simSOX. Its goal is precise: remove the initial sponsor deposit from the optimistic
phase. It should not be confused with a complete redesign of transaction delivery, and it
should not be described as merely ``funding the EntryPoint differently''. At the protocol
level, the important point is that there is no optimistic \(S\)-deposit to reimburse.

This distinction matters because the transaction layer and the protocol economics are
separate. ERC-4337 or a bundler can be used to deliver a transaction, but the
\texttt{no\_S\_deposit} variant is about the economic state of the SOX contract itself. In
the specialized clone implementation, \texttt{OptimisticSOXCloneNoSDeposit} enforces this by
requiring zero value at initialization:

\begin{itemize}
  \item the factory function \texttt{createNoSDeposit} rejects any initial value;
  \item the clone starts in \texttt{WaitPayment}, like the normal path;
  \item the buyer payment required by \texttt{sendPayment} is only the agreed price;
  \item the sponsor tip recorded after payment is set to zero.
\end{itemize}

Thus, simSOX is not a claim that no one pays gas. It is the case where the optimistic part of
the protocol no longer contains an \(S\)-deposit whose purpose is to reimburse external
sponsoring. This makes it a useful comparison point for honest exchanges, because it removes
a piece of accounting that is unnecessary when participants pay their own transaction costs
in the optimistic phase.

The mode is selected before deployment. The frontend and backend expose it as a precontract
variant, and the factory chooses the corresponding implementation. This is important: the
contract does not carry both the normal sponsor-deposit logic and the no-deposit logic at
runtime. The selected clone contains the behavior needed for that mode.

\section[S=B: sponsor is buyer]{\(S=B\): sponsor is buyer}

The \(S=B\) variant is the cleanest precontract-time self-sponsoring case. Since the buyer is
also the optimistic sponsor, the two first steps of the optimistic protocol can be fused in
the implementation. Conceptually, Step 1 is the creation and sponsor deposit by \(S\), while
Step 2 is the buyer payment. If \(S=B\), the same actor is responsible for both actions, so
there is no need to expose a state where a third-party sponsor has created an exchange but
the buyer has not yet committed funds.

The specialized contract \texttt{OptimisticSOXCloneSponsorIsBuyer} implements this fusion.
During initialization, it requires the sponsor to be the buyer and requires enough value to
cover the buyer payment, the completion tip, and the sponsor fees. It then records the buyer
deposit immediately and starts the exchange in state \texttt{WaitKey}. Consequently, the
buyer does not need to call \texttt{sendPayment} as a separate transaction.

This variant addresses the risk identified during the project discussion: if Step 1 and
Step 2 are separated, a sponsor may deploy and fund a contract that is not used because the
buyer never pays. In the \(S=B\) case, the implementation avoids this intermediate state.
The buyer's economic commitment and the creation of the exchange happen together.

The implementation also enforces the role equality on-chain. The factory function
\texttt{createSponsorIsBuyer} requires \texttt{msg.sender} to be the configured buyer, and
the clone itself requires the sponsor address to match the buyer address. This avoids
depending only on frontend selection for a security-relevant role equality.

The \(S=B\) mode is distinct from \texttt{no\_S\_deposit}. In simSOX, the optimistic sponsor
deposit is removed. In \(S=B\), the sponsor role still exists, but it is played by the buyer
and its deposit is provided together with the buyer payment. The two modes therefore answer
different questions and are measured separately.

\section[S=V: sponsor is vendor]{\(S=V\): sponsor is vendor}

The \(S=V\) variant makes the vendor the optimistic sponsor. Like \(S=B\), it is determined
at precontract time and therefore belongs to the optimistic factory selection. Unlike
\(S=B\), however, it does not make the Step 1--2 fusion trivial. The buyer is still a
different party and still has to send the payment after the optimistic contract has been
created.

The dedicated \(S=V\) optimistic clone implements this case. Its initialization requires the
sponsor to be the vendor, requires the sponsor-fee deposit, and then starts in state
\texttt{WaitPayment}. The factory function \texttt{createSponsorIsVendor} also requires the
caller to be the configured vendor. The effect is therefore narrower than in \(S=B\): it
removes the external optimistic sponsor role, but it does not remove the buyer payment
transaction.

This variant was kept because it is a meaningful protocol case and because it was explicitly
requested in the revised scope. It is also useful analytically: it separates the question
``what happens if the optimistic sponsor is not external?'' from the stronger question
``can Step 1 and Step 2 be fused?'' In the current implementation, the answer is yes for
\(S=B\) and no for \(S=V\).

There is a possible future optimization for \(S\ne B\), discussed later in the report: move
the loading of the contract or precontract data to Step 2 and require a signature from
\(S\). Then \(B\) would bring the payment, the precontract data, and \(S\)'s approval in one
transaction. That idea is not part of the implemented \(S=V\) mode because it changes the
deployment architecture more deeply and interacts with ERC-4337 validation.

\section[SB=B and SV=V: self-sponsored disputes]{\(S_B=B\) and \(S_V=V\): self-sponsored disputes}

The equalities \(S_B=B\) and \(S_V=V\) belong to a different category. They are not selected
when the precontract is created. They become relevant only after the buyer indicates that it
is not satisfied and the protocol enters the path that can lead to a dispute.

The first important case is \(S_B=B\). In the paper's step structure, Step 4 is the buyer
announcing dissatisfaction, and Step 5 is the buyer-side dispute sponsor registering and
depositing. If \(S_B=B\), these two actions can be naturally fused: the buyer both expresses
dissatisfaction and supplies the buyer-side dispute sponsor deposit. The clone base provides
\texttt{sendBuyerSelfDisputeSponsorFee} for this self-sponsored case.

For an external \(S_B\), the same fusion is not automatically safe. An external sponsor
should not be able to pretend that the buyer is unhappy. The implementation therefore also
contains an authorization-based entry point for external buyer-side dispute sponsors. The
authorization hash binds the buyer's statement to the chain id, the optimistic contract
address, the buyer, the chosen dispute sponsor, and the commitment. The external \(S_B\)
must provide a signature from the buyer before the contract accepts the deposit. This
directly addresses the concern that a third party could otherwise trigger the buyer-unhappy
path without the buyer's consent.

The second important case is \(S_V=V\). When both \(S_B=B\) and \(S_V=V\) hold, the dispute
does not need the same sponsor-rotation behavior as the fully sponsored case. The losing
party is already the relevant economic party. The implementation therefore selects the
self-sponsored dispute deployer when the buyer dispute sponsor is the buyer and the vendor
dispute sponsor and signer are the vendor. The resulting dispute account uses the
self-sponsored Step 9 logic.

The self-sponsored dispute account specializes Step 9. Instead of relaunching or rotating
the dispute loop, \texttt{handleStep9} moves directly to the appropriate terminal direction:
\texttt{Cancel} when the vendor loses, and \texttt{Complete} otherwise. This is the source of
the expected reduction in full-dispute cost: the optimization is not a small local saving in
one function, but the removal of unnecessary repeated executions of the Step 7--8 loop.

This design also explains why \(S_B=B\) and \(S_V=V\) are not implemented as separate
precontract factory modes. They are dispute-time facts. The same optimistic exchange may
later use either a normal dispute contract or a self-sponsored dispute contract depending on
who actually registers as \(S_B\) and \(S_V\). Treating them as runtime dispute choices keeps
the optimistic contracts smaller and avoids prematurely committing to a dispute path that
may never be used.

\section[Hardcoded SHA-256 circuit variant]{Hardcoded \(\mathsf{SHA256}\) circuit variant}

The hardcoded \(\mathsf{SHA256}\) variant targets a different part of the protocol. It is
not a sponsoring variant. It is a circuit-verification variant for the case where the
description is exactly \(desc=\mathsf{SHA256}(x)\). In the generic circuit mode, the dispute
contract must verify that the gate opened in Step 8 belongs to the committed circuit. This
requires \(h_{\mathit{circuit}}\), the gate bytes, and the Merkle proof usually denoted
\(\pi_1\).

For \(desc=\mathsf{SHA256}\), the circuit structure is determined by public metadata: the
plaintext length, the number of ciphertext blocks, the description hash, and the AES-CTR IV.
The implementation can therefore use a dedicated contract that reconstructs the expected
gate on-chain. This removes the need for the vendor to provide the circuit gate bytes as an
authoritative description and removes the need to verify \(\pi_1\) against
\(h_{\mathit{circuit}}\).

The hardcoded path is fixed at optimistic deployment time through
\texttt{OptimisticSOXAccountHardcodedSHA256}. The constructor stores the hardcoded metadata:
the description hash, the plaintext length, and the ciphertext IV. It also checks that the
number of blocks and the number of gates are consistent with the hardcoded SHA-256 layout.
The choice is therefore not made by the vendor during Step 8. A generic exchange cannot
become hardcoded during a dispute.

The dispute side is handled by \texttt{DisputeSOXAccountHardcodedSHA256} and
\texttt{HardcodedSha256CircuitLib}. In the strict hardcoded tests, the Step 8 call passes
\texttt{gateBytes = 0x} and an empty \texttt{proof1}. The dispute contract ignores these
generic circuit-opening inputs and reconstructs the expected gate bytes itself. It then uses
the same evaluator and accumulator checks needed to verify the values involved in the gate.

This distinction is important for security. The hardcoded contract determines the gate
structure \(g_i\), including the expected opcode and source indices. The vendor may still
provide witness values needed to evaluate the gate, because the smart contract does not
store the entire ciphertext or all intermediate values. These witness values are not trusted
blindly: they are checked against the relevant commitments and Merkle proofs before the gate
result is accepted.

The hardcoded optimization is therefore local to the circuit-membership part of Step 8. It
does not remove the rest of the dispute work. The contract still has to evaluate the
selected gate, verify accumulator proofs for ciphertext or previous values, and execute the
normal dispute state machine. This is why later gas tables distinguish between the cost of a
single Step 8 call and the cost of a complete dispute.

\section{Variants kept, variants discarded, and rationale}

The final implementation keeps the variants that change a meaningful cost component. It
does not implement every theoretically possible role equality as a separate contract. This
is a deliberate design decision rather than an omission.

\begin{table}[htbp]
  \centering
  \begin{tabularx}{\linewidth}{>{\raggedright\arraybackslash}p{0.25\linewidth}>{\raggedright\arraybackslash}p{0.22\linewidth}>{\raggedright\arraybackslash}X}
    \toprule
    Case & Status & Rationale \\
    \midrule
    Normal sponsored SOX &
    Implemented &
    Required baseline for comparison with the inherited sponsored architecture. \\
    \texttt{no\_S\_deposit} / simSOX &
    Implemented &
    Removes the optimistic sponsor deposit and isolates the cost of the simplified honest
    path. \\
    \(S=B\) &
    Implemented &
    Gives a concrete and safe Step 1--2 fusion because the buyer and optimistic sponsor are
    the same actor. \\
    \(S=V\) &
    Implemented &
    Removes the external optimistic sponsor role while keeping the buyer payment as a
    separate action. \\
    \(S_B=B\) and \(S_V=V\) &
    Implemented as dispute-time specialization &
    Avoids unnecessary Step 9 repetitions when both dispute sponsors are the corresponding
    parties. \\
    Hardcoded \(\mathsf{SHA256}\) &
    Implemented as a dedicated path &
    Removes \(h_{\mathit{circuit}}\) and \(\pi_1\) from the dedicated SHA-256 Step 8
    verification path. \\
    All 12 role-combination cases as separate optimistic contracts &
    Not implemented separately &
    Several combinations differ only after Step 4 or only in accounting labels. Encoding all
    combinations as optimistic contracts would increase code size without a proportional
    measurement benefit. \\
    Cases such as \(S_B=V\) or \(S_V=B\) &
    Documented but not specialized &
    These edge cases are possible role equalities, but they were not expected to change the
    main gas evaluation in the retained scenarios. \\
    Arbitrary hardcoded descriptions &
    Not implemented &
    The project only implements the dedicated \(desc=\mathsf{SHA256}\) layout. Other
    descriptions would need their own circuit-specific verifier. \\
    \bottomrule
  \end{tabularx}
  \caption{Retained and non-retained variant cases.}
  \label{tab:variant-rationale}
\end{table}

The resulting architecture has three levels of specialization. First, the frontend and
backend choose a precontract variant among normal, simSOX, \(S=B\), and \(S=V\). Second, the
factory deploys the corresponding minimal optimistic clone, so unused optimistic branches
are not carried into every exchange. Third, if a dispute occurs, the optimistic contract
selects the normal or self-sponsored dispute deployer based on the actual dispute sponsors.
The hardcoded SHA-256 path is kept as a dedicated circuit-specialized path because its
optimization concerns Step 8 rather than the optimistic sponsor logic.

This organization is the foundation for the implementation architecture in the next
chapter. The important point is that the final design does not try to make a single smart
contract smarter. It tries to make each deployed contract narrower. In a gas-sensitive
protocol, that distinction is essential.
